\chapter{The dimensional transforms}

\textbf{Purpose:} Record what each transform in this work translates between, what the translation buys, and the single condition under which any of them helps at all. \textbf{Scope:} the bitpack, fixed point, the number-theoretic transform, the four-step fold, the golden placement, the harmonic expansion, and the permutation read as an orthogonal map.

Seven transforms carry this work. They look unrelated and they do one thing: they express the same information in coordinates where a question is easier to ask. No transform here creates information or destroys it. Each has an inverse or is lossless at full rank.

\section{The seven}

\textbf{The bitpack.} An integer and a run of 32-bit limbs are the same bytes named twice. \texttt{to\_bytes} and \texttt{from\_bytes} are not conversions; they are two vocabularies for one object. This lets a number be handed to a device that only knows arrays, without anything being computed.

\textbf{Fixed point.} A real to whatever depth is paid for, held as an integer scaled by a power of two. The reals are the human coordinate and the integer is the machine's, and the translation is exact in both directions. The whole floor question turns on this: a residual stopping at \(10^{-16}\) is speaking float64, and one falling with the width is speaking the object.

\textbf{The number-theoretic transform.} Coefficients to evaluations and back. A polynomial of \(n\) coefficients evaluated at \(n\) points costs \(n^2\) pointwise and \(n \log n\) through the transform, because the transform shares work between the points that pointwise evaluation repeats.

\textbf{The four-step fold.} One dimension of length \(N\) becomes two of \(P \times Q\), with a twiddle between them. Nothing is lost, since it is a relabeling of the same indices, and it lets a transform too large for one pass fit in two.

\textbf{The golden placement.} An index becomes a direction. Index \(k\) goes to height \(1 - 2(k + 0.5)/n\) at longitude \(k\gamma\), and \(256\) bit positions become \(256\) points on a sphere with no seam and no pole pile. A one-dimensional list becomes a two-dimensional surface.

\textbf{The harmonic expansion.} A set of points on that surface becomes \((L+1)^2\) coefficients sorted by angular scale. At full rank it is lossless, making it a change of vocabulary and not a summary, and the sorting is what it buys: structure at a scale concentrates in one band instead of smearing across \(256\) numbers.

\textbf{The permutation as an orthogonal map.} \(k\) observations of a quantity become a vector in \(k\) dimensions, and correlation becomes the cosine of an angle. A permutation is then a rotation within the coordinate subgroup, preserving every length and randomizing every angle. This one was hiding in plain sight all night as \emph{shuffling the data}.

\section{What they preserve and what they move}

Every one preserves the information and moves the coordinates. That family resemblance has a consequence worth stating flatly: \textbf{a transform cannot make a question answerable that was unanswerable. It can only make an answerable question cheap.}

The engine is faster because of the NTT, and not more capable. The boundary reading is legible because of the harmonic expansion, and not more informed. Nothing in this tree computes anything a slower arrangement could not compute, and stating that plainly keeps the rest of the claims defensible.

\section{The single condition: rank deficiency}

A transform pays when the object has fewer coefficients than it has points. That is the entire mechanism, in every case.

Multipoint evaluation is fast because a degree-\(n\) polynomial has \(n\) coefficients however many places it is evaluated. The work at one point overlaps the work at the next. Harmonic sorting is useful because a smooth field has most of its power in low degrees, and the high ones can be dropped. The arms pool because independent observations share a mean, and \(k\) of them carry one number's worth of uncertainty instead of \(k\).

That fixes when a transform is useless, and the test for it is direct. A Boolean function of degree below \(k\) sums to zero over any \(k\)-dimensional affine subspace, and a cube of inputs therefore collapses if and only if the function is rank-deficient in that sense. Measured over \(126\) cells at every depth from eight rounds to sixty-four, nothing collapsed: every cell sat near \(128\) of \(256\), which is full degree.

\textbf{An object at full rank in a basis has as many coefficients as points, and no change of basis helps.} Not \emph{the transform is slow}: there is nothing for it to share. No transform in this tree could therefore shortcut a hash, and every one of them speeds up arithmetic, because polynomials are rank-deficient by construction and a hash is built specifically not to be.

\section{Where a translation earned its keep, and where it did not}

\textbf{Earned it.} The permutation-as-rotation reading explained why independent observations discriminate so sharply, and did it better than the counting argument it replaced: in \(k\) dimensions random vectors are nearly perpendicular, with typical cosine \(\sqrt{2/\pi k}\), giving \(0.051\) at \(245\) dimensions and \(0.025\) at \(1000\). The discrimination is concentration of measure, and not sample size, and that framing predicts the right scaling where counting did not.

\textbf{Earned it.} The harmonic expansion beat every rival scanner on spatial structure by a factor of twenty-two, \(328\) standard errors against \(14.8\), because sorting by angular scale gives the coordinate in which a shape is one coefficient instead of \(256\).

\textbf{Did not.} The same expansion read \(2.43\) against its own \(3.75\) floor on a forced parity, below its own noise. A parity is not a linear property of which points are lit. The expansion is blind to it by construction, and the crudest scan in the comparison caught it. A transform translates into one language and is deaf in every other.

\textbf{Did not.} Reading a state through a mirrored placement looked like a chirality test and was a geometric identity: reflection negates every longitude, conjugating every coefficient and negating every phase, for any data whatever. The answer was fixed before the data was consulted. A transform's own symmetries can produce a result saying nothing about the object.

\section{The practical form}

\begin{enumerate}
  \item Ask what coordinate the question is naturally cheap in.
  \item Check the object is rank-deficient in that coordinate, because otherwise the move is free of benefit as well as free of cost.
  \item Apply the transform, which changes nothing but the vocabulary.
  \item Check the result is not a symmetry of the transform instead of a property of the object. Run the same reading on data with nothing in it and see whether the answer moves.
  \item Translate back, and state the finding in the coordinate the question was asked in.
\end{enumerate}

Step four is what this work learned the hard way, twice.
